\documentclass[12pt,a4paper]{article}

% ── Packages ────────────────────────────────────────────────────────────────
\usepackage[a4paper, margin=2.5cm]{geometry}
\usepackage{graphicx}
\usepackage{hyperref}
\usepackage{hypcap}
\usepackage{bookmark}
\usepackage{amsmath,amssymb,amsfonts}
\usepackage{titlesec}
\usepackage{caption}
\usepackage{subcaption}
\usepackage{enumitem}
\usepackage{lmodern}
\usepackage{cmap}
\usepackage{xcolor}
\usepackage{tikz}
\usetikzlibrary{shapes.geometric, arrows, positioning}
\usepackage{float}
\usepackage{colortbl}


% Soft pastel colors for headings
\definecolor{softblue}{HTML}{6B9AC4}
\definecolor{softteal}{HTML}{7FB3B3}
\definecolor{softpurple}{HTML}{B39EB5}
\definecolor{low}{HTML}{F8E6E6}
\definecolor{mid}{HTML}{E6F2F2}
\definecolor{high}{HTML}{D6EAF8}
\definecolor{best}{HTML}{A9CCE3}
% Table and long table support for large result tables
\usepackage{tabularx}
\usepackage{longtable}
\usepackage{booktabs}
\usepackage{multirow}
\usepackage{array}
\usepackage{fancyhdr}
\usepackage[T1]{fontenc}
\usepackage[utf8]{inputenc}
\usepackage{comment}
\usepackage{lmodern}
\usepackage{cmap}
\usepackage{chngcntr}
\usepackage{placeins}
\usepackage{tcolorbox}
\counterwithin{figure}{section}
\counterwithin{table}{section}


% ── Hyperref setup ───────────────────────────────────────────────────────────
\hypersetup{
	colorlinks=true,
	linkcolor=black,
	urlcolor=black,
	citecolor=black
}

% ── Section formatting (matches PDF style) ───────────────────────────────────
\titleformat{\section}{\large\bfseries\color{softblue}}{\thesection}{1em}{}
\titleformat{\subsection}{\normalsize\bfseries}{\thesubsection}{1em}{}
\titleformat{\subsubsection}{\normalsize\bfseries}{\thesubsubsection}{1em}{}

% ── Header / Footer ──────────────────────────────────────────────────────────
\pagestyle{fancy}
\fancyhf{}
\rhead{\small Neural Networks -- Final Project MCQ}
\lhead{\small Cairo University}
\cfoot{\thepage}
\setlength{\headheight}{14pt}
\renewcommand{\headrulewidth}{0.4pt}

% ── Graphics path ────────────────────────────────────────────────────────────
\graphicspath{
    {Figures/}
    {../prob_1/Figures/}
	{../prob_2/Figures/}
	{../prob_3/Figures/}
	{../prob_4/Figures/}
	{../prob_5/Figures/}
	{../prob_6/Figures/}
}
\begin{document}

% ── Title Page ───────────────────────────────────────────────────────────────
\begin{titlepage}
\centering
\vspace*{0.5cm}

\makebox[\textwidth][c]{%
\raisebox{-0.4cm}{\includegraphics[width=0.25\textwidth]{Figures/image1.png}}\hspace{2.5cm}%
\includegraphics[width=0.17\textwidth]{Figures/image2.png}%
}

\vspace{0.5cm}
{\Large\bfseries Cairo University\\}
\vspace{0.2cm}
{\large Faculty of Engineering\\}
\vspace{0.2cm}
{\normalsize Electronics and Electrical Communications Engineering\\}

\vspace{1.2cm}
\rule{0.8\textwidth}{1.5pt}\\[0.4cm]
{\LARGE\bfseries Final Project MCQ}\\[0.3cm]
{\large\bfseries ELC 4028 -- Artificial Neural Networks}\\[0.2cm]
\rule{0.8\textwidth}{1.5pt}

\vspace{1.2cm}

\begin{tabular}{|c|c|l|}
\hline
\textbf{ID} & \textbf{Section} & \textbf{Name} \\
\hline
9220399 & 2 & Shahd Gamal Mahmoud Hussein \\
\hline
9220411 & 2 & Salah El-Din Hassan Hassan Mohamed \\
\hline
9220984 & 4 & Yousef Khaled Omar Mahmoud \\
\hline
\end{tabular}

\vspace{1.2cm}

\textit{Supervised by:}\\[0.3cm]
\textit{Dr.\ Mohsen Rashwan}\\
\textit{Dr.\ Mohamed Motaz}

\vfill
\end{titlepage}



% =====================================================
% Project Workflow Diagram
% =====================================================

\definecolor{softblue}{RGB}{173,216,230}
\definecolor{softgreen}{RGB}{190,230,200}
\definecolor{softpurple}{RGB}{220,210,255}
\definecolor{softorange}{RGB}{255,220,180}
\definecolor{softred}{RGB}{255,200,200}

\subsection*{Q1. What does RAG stand for?}

\begin{enumerate}[label=\Alph*.]
\item Retrieval-Augmented Generation
\item Random Attention Generator
\item Retrieval Artificial Graph
\item Recursive Augmented Graph
\end{enumerate}

\begin{tcolorbox}[colback=mid,colframe=softblue,title=Correct Answer]
A. Retrieval-Augmented Generation

RAG combines information retrieval systems with language generation models to improve answer quality and reduce hallucinations.
\end{tcolorbox}


\subsection*{Q2. What is the main purpose of RAG systems?}

\begin{enumerate}[label=\Alph*.]
\item Compress neural networks
\item Retrieve external knowledge before generation
\item Replace transformers completely
\item Reduce dataset size
\end{enumerate}

\begin{tcolorbox}[colback=mid,colframe=softblue,title=Correct Answer]
B. Retrieve external knowledge before generation

RAG systems first retrieve relevant documents and then use them to generate context-aware responses.
\end{tcolorbox}


\subsection*{Q3. Which problem does RAG primarily help reduce?}

\begin{enumerate}[label=\Alph*.]
\item Overfitting
\item Vanishing gradients
\item Hallucination
\item Image noise
\end{enumerate}

\begin{tcolorbox}[colback=mid,colframe=softblue,title=Correct Answer]
C. Hallucination

RAG reduces hallucinations by grounding responses using retrieved external information.
\end{tcolorbox}


\subsection*{Q4. What are embeddings in RAG systems?}

\begin{enumerate}[label=\Alph*.]
\item Image compression algorithms
\item Numerical vector representations of text
\item Audio processing layers
\item Database encryption methods
\end{enumerate}

\begin{tcolorbox}[colback=mid,colframe=softblue,title=Correct Answer]
B. Numerical vector representations of text

Embeddings convert text into vectors that capture semantic meaning and relationships.
\end{tcolorbox}


\subsection*{Q5. Which database type is commonly used in RAG systems?}

\begin{enumerate}[label=\Alph*.]
\item Relational database
\item Vector database
\item Blockchain database
\item Graphical database
\end{enumerate}

\begin{tcolorbox}[colback=mid,colframe=softblue,title=Correct Answer]
B. Vector database

Vector databases store embeddings and perform semantic similarity search efficiently.
\end{tcolorbox}


\subsection*{Q6. Which similarity metric is commonly used for retrieval?}

\begin{enumerate}[label=\Alph*.]
\item Euclidean sorting
\item Binary matching
\item Cosine similarity
\item Histogram equalization
\end{enumerate}

\begin{tcolorbox}[colback=mid,colframe=softblue,title=Correct Answer]
C. Cosine similarity

Cosine similarity measures the angular similarity between embedding vectors.
\end{tcolorbox}


\subsection*{Q7. What happens during the retrieval phase in RAG?}

\begin{enumerate}[label=\Alph*.]
\item Model weights are retrained
\item Relevant document chunks are searched
\item Images are generated
\item Audio signals are filtered
\end{enumerate}

\begin{tcolorbox}[colback=mid,colframe=softblue,title=Correct Answer]
B. Relevant document chunks are searched

The retrieval stage identifies semantically similar documents related to the user query.
\end{tcolorbox}


\subsection*{Q8. What is chunking in RAG systems?}

\begin{enumerate}[label=\Alph*.]
\item Compressing neural weights
\item Splitting documents into smaller sections
\item Removing embeddings
\item Encrypting text files
\end{enumerate}

\begin{tcolorbox}[colback=mid,colframe=softblue,title=Correct Answer]
B. Splitting documents into smaller sections

Large documents are divided into smaller chunks to improve retrieval accuracy.
\end{tcolorbox}


\subsection*{Q9. Which of the following is an advantage of RAG compared to fine-tuning?}

\begin{enumerate}[label=\Alph*.]
\item Requires retraining for every update
\item Easier knowledge updates
\item Uses no embeddings
\item Removes transformers completely
\end{enumerate}

\begin{tcolorbox}[colback=mid,colframe=softblue,title=Correct Answer]
B. Easier knowledge updates

RAG allows updating external documents without retraining the entire model.
\end{tcolorbox}


\subsection*{Q10. Which ANN-related concept is strongly connected to RAG?}

\begin{enumerate}[label=\Alph*.]
\item Backpropagation only
\item Embeddings and transformers
\item Edge detection only
\item Fourier transforms
\end{enumerate}

\begin{tcolorbox}[colback=mid,colframe=softblue,title=Correct Answer]
B. Embeddings and transformers

RAG relies heavily on transformer-based models and semantic vector embeddings.
\end{tcolorbox}


\subsection*{Q11. What is semantic search?}

\begin{enumerate}[label=\Alph*.]
\item Searching by exact keyword only
\item Searching based on meaning and context
\item Searching image colors
\item Searching file size
\end{enumerate}

\begin{tcolorbox}[colback=mid,colframe=softblue,title=Correct Answer]
B. Searching based on meaning and context

Semantic search retrieves information using contextual meaning rather than exact keyword matching.
\end{tcolorbox}


\subsection*{Q12. Why are vector databases important in RAG systems?}

\begin{enumerate}[label=\Alph*.]
\item They store operating systems
\item They perform fast similarity search
\item They replace embeddings
\item They generate neural networks
\end{enumerate}

\begin{tcolorbox}[colback=mid,colframe=softblue,title=Correct Answer]
B. They perform fast similarity search

Vector databases efficiently retrieve semantically similar embeddings during the retrieval phase.
\end{tcolorbox}

\subsection*{Q13. What is the role of the embedding model in RAG systems?}

\begin{enumerate}[label=\Alph*.]
\item Generate images from text
\item Convert text into vector representations
\item Compress the database
\item Remove irrelevant documents
\end{enumerate}

\begin{tcolorbox}[colback=mid,colframe=softblue,title=Correct Answer]
B. Convert text into vector representations

Embedding models transform text into numerical vectors that capture semantic meaning for similarity search.
\end{tcolorbox}


\subsection*{Q14. Which component generates the final response in a RAG pipeline?}

\begin{enumerate}[label=\Alph*.]
\item Vector database
\item Embedding layer
\item Large Language Model (LLM)
\item Data loader
\end{enumerate}

\begin{tcolorbox}[colback=mid,colframe=softblue,title=Correct Answer]
C. Large Language Model (LLM)

After retrieval, the LLM uses the retrieved context to generate the final answer.
\end{tcolorbox}


\subsection*{Q15. What is one limitation of RAG systems?}

\begin{enumerate}[label=\Alph*.]
\item They cannot use transformers
\item Poor retrieval can produce incorrect answers
\item They completely eliminate latency
\item They require no storage
\end{enumerate}

\begin{tcolorbox}[colback=mid,colframe=softblue,title=Correct Answer]
B. Poor retrieval can produce incorrect answers

If irrelevant or incorrect chunks are retrieved, the generated response quality decreases significantly.
\end{tcolorbox}


\subsection*{Q16. Which of the following tools is commonly used for vector similarity search?}

\begin{enumerate}[label=\Alph*.]
\item FAISS
\item Photoshop
\item AutoCAD
\item MySQL Workbench
\end{enumerate}

\begin{tcolorbox}[colback=mid,colframe=softblue,title=Correct Answer]
A. FAISS

FAISS is a popular library developed for efficient vector similarity search and retrieval.
\end{tcolorbox}


\subsection*{Q17. Why is chunk size important in RAG systems?}

\begin{enumerate}[label=\Alph*.]
\item It controls internet speed
\item It affects retrieval quality and context understanding
\item It changes GPU temperature
\item It removes embeddings
\end{enumerate}

\begin{tcolorbox}[colback=mid,colframe=softblue,title=Correct Answer]
B. It affects retrieval quality and context understanding

Very small chunks may lose context, while very large chunks may reduce retrieval precision.
\end{tcolorbox}


\subsection*{Q18. Which type of search is performed in vector databases?}

\begin{enumerate}[label=\Alph*.]
\item Lexical search only
\item Semantic similarity search
\item Binary search only
\item Pixel-based search
\end{enumerate}

\begin{tcolorbox}[colback=mid,colframe=softblue,title=Correct Answer]
B. Semantic similarity search

Vector databases retrieve information based on semantic meaning rather than exact word matching.
\end{tcolorbox}


\subsection*{Q19. What is prompt augmentation in RAG systems?}

\begin{enumerate}[label=\Alph*.]
\item Reducing prompt size
\item Adding retrieved context to the prompt
\item Encrypting the prompt
\item Removing user queries
\end{enumerate}

\begin{tcolorbox}[colback=mid,colframe=softblue,title=Correct Answer]
B. Adding retrieved context to the prompt

Retrieved documents are appended to the user query to improve answer accuracy.
\end{tcolorbox}


\subsection*{Q20. Why are RAG systems considered useful for modern AI applications?}

\begin{enumerate}[label=\Alph*.]
\item They allow AI systems to use updated external knowledge
\item They completely replace neural networks
\item They remove the need for databases
\item They eliminate all AI errors
\end{enumerate}

\begin{tcolorbox}[colback=mid,colframe=softblue,title=Correct Answer]
A. They allow AI systems to use updated external knowledge

RAG systems improve reliability by integrating external and dynamically updated information sources.
\end{tcolorbox}


\end{document}